% =============================================================================
% 5 REFLECTION AND CONCLUSION        budget: 1.00 page (10-15 % of the report)
% -----------------------------------------------------------------------------
% LITERATURE: ALMOST NONE  ->  0-2 citations, and no new sources
%   A conclusion that introduces fresh literature reads as if the argument was
%   not finished in the main part. This chapter is your own voice: it judges
%   what the earlier chapters did. Reflection cannot be delegated to a source.
%   The one or two places a citation is defensible:
%     - when naming a model that was adapted, point back to the source already
%       used for it in the main part (same key, same page) so the reader can
%       see what was changed against what
%     - the closing claim that there is no single best methodology and that
%       hybrids are legitimate. Wysocki is the obvious source; an empirical
%       paper on agile adoption failures would be stronger and is available -
%       the "evidence__*" PDFs in literature/ are that kind of source. This
%       single citation stops the ending from sounding like an opinion.
%   Do NOT: pad this chapter with citations to look academic. Everything here
%   should be traceable to YOUR chapters 1-4, not to new reading.
%
% WHAT YOU WANT TO SAY HERE, IN SHORT
%   - "The project produced a scoped product, ten stories, a prioritised
%      backlog and a planned first sprint - measured against the objective in
%      chapter 1, this is what was and was not achieved." [own]
%   - "These models were simplified or dropped, and for these reasons: MoSCoW
%      over WSJF, no acceptance criteria, product scope over project scope,
%      the chosen sprint length and cadence, estimation kept minimal." [own -
%      this paragraph is where Transfer is scored]
%   - "Time, budget and people: this is how the approach used them, and this is
%      where it would strain." [own]
%   - "Done again, this would change." [own]
%   - "For further study and professional practice, this follows." [own]
%   - "No methodology is universally right; agile fits this product, hybrids
%      are legitimate elsewhere." [own + the one closing citation]
% -----------------------------------------------------------------------------
% THIS CHAPTER IS MANDATORY AND MUST NOT REPEAT CHAPTER 4. Both tutors say so
% independently:
%   "The final task is to make a comparison. But in addition, every written
%    assignment should end with a reflection that looks at the whole project.
%    So there should be an additional conclusion" [SCH, 17.12.22]
%   Report = introduction, body, CONCLUSION, bibliography [G-Jan26].
% [TASK] adds: "Document your process in detail and reflect on it continuously,
% but at the latest in the final chapter."
%
% WHY IT IS WORTH THE PAGE: Transfer (15 %) is defined as "appropriate transfer
% of theories/models, conclusions in the final remarks" - it is scored here
% almost entirely. Process (25 %) and Quality (25 %) are also decided by the
% depth of this reflection. Do not economise.
%
% [GL] CONCLUDING KEY QUESTIONS:
%   - What conclusions can be drawn from the results and findings presented in
%     the main section?
%   - How do these findings relate to your future studies as well as personal
%     and professional development?
%
% -----------------------------------------------------------------------------
% TRANSFER - the highest-yield paragraph of the chapter
% [TASK], verbatim: "You may have solidified or modified certain steps. You
% should describe and explain this." So name every place where a model was
% simplified, adapted or rejected, and give the reason:
%   - MoSCoW chosen, WSJF examined and rejected as too pseudo-scientific for a
%     ten-item backlog
%   - acceptance criteria deliberately omitted, because the stories stay at the
%     conversation stage in sprint zero
%   - the product scope defined in full, the project scope only sketched
%   - the chosen sprint length and refinement cadence, and what was traded away
%   - estimation kept at the level of sprint planning, without planning poker
%   - the first sprint deliberately underloaded because a new team estimates
%     its own capacity badly
%   - Scrum applied as described in the Scrum Guide, but with a fictitious
%     team of a stated size - say where reality would differ
%
% RESOURCES (10 %) - explicitly reflect on time, budget and people. [TASK]:
% "Make sure that in the projects, resources such as time, budget and people
% are used efficiently." Two lines of leverage already in the report: the
% inverted iron triangle from chapter 4, and the PO's "here's the impact if you
% want this in" from 3.1.
%
% PROCESS (25 %) - state what worked in the approach and what would be done
% differently. [GL] key words: goal achievement (was the objective met, and if
% not why), effects and changes achieved, areas requiring improvement,
% continuous improvements to implement, how sustainable structures could be
% created to anchor the results.
%
% PERSONAL AND PROFESSIONAL DEVELOPMENT - answer the [GL] question directly:
% what this planning exercise implies for further studies and for the writer's
% own practice. One short paragraph, still in the third person.
%
% -----------------------------------------------------------------------------
% CLOSING TONE - the safe ending, in his own words [G-Feb25]:
%   "there is no perfect methodology for systems development"
%   hybrids of waterfall and agile are "quite possible to do"
% Do NOT end on "agile is simply better". A balanced close is both more
% defensible and better aligned with the Quality criterion.
%
% -----------------------------------------------------------------------------
% FINAL PASS BEFORE SUBMISSION (see CHECKLIST.md for the full list)
%   - 7-10 pages of text, tables included, appendix absent
%   - all three task questions balanced, none noticeably thinner
%   - every citation carries a page number; every source in the list is cited
%     in the text
%   - no first person anywhere, one consistent tense
%   - no figures; only Tables 1 and 2
%   - list of abbreviations trimmed to what is actually used
%   - anti-plagiarism pledge submitted BEFORE uploading
% =============================================================================

% #############################################################################
% WRITING PLAN - CHAPTER 5
% Budget: 1.00 page = ~700 WORDS across 6 paragraphs => ~117 WORDS EACH.
% 0-2 CITATIONS, AND NO NEW SOURCES.
%
% ** THE RULE THAT SHAPES EVERYTHING HERE: ** a conclusion that introduces fresh
% literature reads as if the argument was not finished in the main part. Only
% two kinds of citation are legitimate:
%   (a) pointing BACK to a source already used, to show what was changed against
%       what;
%   (b) the single closing "no perfect methodology" claim.
%
% This chapter is where Transfer (15 %) is scored almost entirely, and Process
% (25 %) and Quality (25 %) are decided by the depth of the reflection. It is
% also the chapter students most often thin out. Do not economise.
%
% !! IT MUST NOT REPEAT CHAPTER 4. Chapter 4 compares the two approaches;
%    chapter 5 judges what THIS project did. Both tutors say this independently.
%
% -----------------------------------------------------------------------------
% ONE SOURCE WAS DELIBERATELY MOVED OUT OF THIS CHAPTER. Suryaatmaja et al.
% (2020) was originally planned for chapter 5 only - but a source used ONLY in
% the conclusion is a new source in the conclusion by construction. It now
% appears in 3.1 (the PO role undefined in practice), which makes a back-
% reference here legitimate under rule (a).
% -----------------------------------------------------------------------------
% !! THE MANIFESTO LOCATOR FOOTNOTE. If principle 10 is used below and no Beck
%    citation has appeared earlier, the footnote goes here. Exact wording:
%
%    "The Agile Manifesto and its accompanying principles are published only as
%     unpaginated web pages, so APA 7 requires them to be cited by paragraph
%     rather than page; the four values and twelve principles are presented
%     unnumbered in the source, and the numbering used here is the present
%     author's, counting from the first principle."
%
%    This matters more than it looks: the module's one named point-loser is a
%    missing page number, so explaining why THIS source legitimately has none
%    converts a potential defect into a visible display of source care.
% #############################################################################

\section{Reflexion und Fazit}

% --- P1: what the process produced, measured against the objective -------
% NO CITATION. Own voice throughout.
% Say: measured against the objective stated in chapter 1, the project produced
% a scoped product, ten user stories, a prioritised initial backlog and a
% planned first sprint. State plainly what was NOT produced - no increment, no
% velocity, no multi-sprint history - and why that is the correct end point
% rather than a shortfall: the task is a planning exercise and the report ends
% at sprint zero.
% [GL] also expects the main part to record "milestones and any difficulties" -
% if chapter 3 did not carry a sentence on what refinement actually changed and
% what was hard, carry it here.
Gemessen an dem in Kapitel~1 gesetzten Ziel hat das Projekt einen abgegrenzten
Produktscope für Cookbox, zehn User Stories in einem einheitlichen Format, ein
priorisiertes initiales Backlog aus diesen Stories und einen geplanten ersten
Sprint hervorgebracht. Es hat keine lauffähige Software, keine Velocity und
keine Sprint-Historie hervorgebracht, was für eine Planungsübung, die bei Sprint
Zero endet, der richtige Endpunkt und kein Defizit ist. Ein Ergebnis verdient es,
klar benannt zu werden: Die ursprüngliche Reihenfolge war innerhalb eines
einzigen Sprints zweimal falsch. US-08 war zu niedrig und US-07 zu hoch
eingestuft, und das Refinement hat beides aufgedeckt. Das ist ein Argument für
den Prozess und kein Beleg dafür, dass das erste Backlog richtig war.

Mehrere Modelle wurden in angepasster Form angewendet, und die Prinzipien hinter
dem Agilen Manifest liefern den Maßstab, an dem die Auslassungen zu lesen sind:
\enquote{Simplicity-{}-the art of maximizing the amount of work not done-{}-is
essential} \parencite[para.~10]{beck2001principles}. MoSCoW wurde unverändert
übernommen; Weighted Shortest Job First wurde geprüft und verworfen, weil sich
bei einem Backlog aus zehn Einträgen und einem einzigen Product Owner eine
Bewertungsformel nicht lohnt, deren Eingangsgrößen selbst Schätzungen sind. Das
Backlog wurde dahingehend abgewandelt, dass seine geordneten Einträge auf User
Stories beschränkt wurden, was Cohns dokumentierte Änderung am reinen Scrum ist
und nicht Scrum selbst. Akzeptanzkriterien, der dritte Bestandteil des in
Kapitel~2 beschriebenen Konzepts der User Story, wurden weggelassen, weil sie in
den Sprint gehören, in dem eine Story gebaut wird. Der Produktscope wurde
vollständig definiert, der Projektscope nur skizziert. Die Refinement-Taktung von
zwei einstündigen Sitzungen je Sprint wurde aus veröffentlichten Praxiswerten und
einer veröffentlichten Zeitobergrenze abgeleitet statt einer Norm entnommen und
auf ein Team zugeschnitten, das die Domäne neu kennenlernt; ein erfahrenes Team
käme mit weniger aus. Die Schätzung beschränkt sich auf das, was Sprint Planning
verlangt, da hier das Refinement und nicht der Aufwand die Frage ist, und der
erste Sprint wurde mit zwei der vier Must-haves bewusst unterladen. Im Übrigen
wurde Scrum wie beschrieben angewendet, allerdings auf ein Team, das nur auf dem
Papier existiert: Ein reales Team hätte die Schätzung, die US-07 zurückstufte,
nicht so schnell erreicht, wie dieser Bericht unterstellt.

Zeit, Geld und Personal stehen in diesem Projekt allesamt fest --- sechs Monate,
240.000~EUR und ein Team aus sechs Personen ---, wodurch der Scope als einzige
Variable bleibt, und das ist zuerst eine Ressourcen- und dann eine
Planungsentscheidung. Sie ist der Grund, warum das Moderationswerkzeug von US-09
gestrichen wurde und das Support-Team die Startsammlung zum Launch von Hand
moderiert, statt den Termin zu verschieben. Der charakteristische Zug des Product
Owners folgt derselben Rechnung: Ein später Wunsch wird nicht abgelehnt, sondern
mit der Arbeit bepreist, die er verdrängt. Derselben Disziplin unterliegt dieser
Bericht, denn die Seitenbegrenzung ist selbst eine ausgewiesene
Ressourcenanforderung.

Ein zweites Mal würde die Reihenfolge mit den Entwicklern im Raum festgelegt und
die Rechtslage zuerst geprüft. Zwei der drei in Abschnitt~\ref{subsec:refinement}
festgehaltenen Umsortierungen beruhten auf Informationen, die es schon vor dem
Festlegen des Backlogs gab: Die Entwickler hätten die Einheitenumrechnung hinter
US-07 in der ersten Woche bepreisen können, und der rechtliche Status der
Kontolöschung war Aktenlage und keine Ermessensfrage. Allein nach Wert und
Abhängigkeit zu ordnen und das Team erst danach hinzuzuziehen, hat zwei der
Refinement-Entscheidungen des ersten Sprints vermeidbar gemacht. Was sich nicht
ändern würde, ist der Rhythmus: Das Refinement war ein fester Termin und keine
Reaktion, weshalb beide Fehler als gewöhnliche Korrekturen auftraten und nicht
als Neuplanung.

Für das weitere Studium verortet die Übung die Schwierigkeit agiler Planung
genau. Zehn Stories zu schreiben geht schnell; ihre Reihenfolge vertretbar zu
halten, wenn neue Informationen eintreffen, ist der Teil, der Übung verlangt, und
darauf bauen spätere Arbeiten zu Schätzung und Auslieferung auf. Für die
berufliche Praxis übertragen sich drei Gewohnheiten über Software hinaus: zu
benennen, wem eine Anforderung dient und was sie wert ist, bevor sie eingeordnet
wird, wozu der Nutzen-Teilsatz jeder Cookbox-Story gezwungen hat; die Menschen
schätzen zu lassen, die die Arbeit tun werden; und jede Arbeitsreihenfolge als
verfallend zu behandeln. Die Artefakte waren billig, und ein großer Teil ihres
Werts stammte aus den Sitzungen, die sie überarbeitet haben.

Ob das Vorgehen zu diesem Produkt gepasst hat, ist eine engere Frage als die, ob
Agilität besser ist. Zu Cookbox passt es: Diejenige Story, die eine schriftliche
Beschreibung am ehesten verfehlt --- die Kochansicht aus US-03 ---, ist jene, die
früh auf einem echten Gerät gesehen werden musste. Ein Produkt mit einer
festgeschriebenen regulatorischen Spezifikation oder ein Vertrag mit
Bezahlung je Liefergegenstand würde eine andere Antwort rechtfertigen, und ein
hybrides Vorgehen bleibt legitim. Das Project Management Institute formuliert den
allgemeinen Fall:
\enquote{There is no single approach that can be applied to all projects all of
the time} \parencite[PMBOK Guide, Section~3.2]{pmbok2021}.


% --- P2: TRANSFER - models adapted, simplified or rejected --------------
% ** THE HIGHEST-YIELD PARAGRAPH IN THE CHAPTER. ** The task PDF says it
% verbatim: "You may have solidified or modified certain steps. You should
% describe and explain this." So NAME each place a model was changed, and give
% the reason. Mostly own voice; at most ONE back-reference citation.
%
% ** USE PMBOK'S OWN TAILORING VOCABULARY AS THE SCAFFOLD, UNCITED: **
%   add / modify / remove / blend / align (PMBOK Guide, Section 3.3.2).
%   Naming which of the five each decision was is the cheapest way to make the
%   paragraph read as competence rather than as a list of shortcuts. Using the
%   vocabulary without quoting it costs no citation.
%
% The list to work through (each needs a REASON, not just a mention):
%   - MoSCoW ADOPTED, WSJF EXAMINED AND REJECTED - and say on what grounds:
%     fit to a ten-item backlog with one PO, not "value cannot be quantified".
%     (Chapter 3.2 made this argument; here, judge it.)
%   - THE BACKLOG CONSTRAINED TO USER STORIES ONLY. Flag that this is Cohn's
%     documented modification of plain Scrum, not Scrum itself - chapter 3.2
%     already cited him for it, so a back-reference is legitimate here if the
%     citation budget allows.
%   - ACCEPTANCE CRITERIA DELIBERATELY OMITTED. ** This is the best available
%     use of the (a) back-reference rule: **
%       \parencite[p.~383]{lucassen2016}
%       The concept of a user story formally includes acceptance criteria as its
%       third component; this report does not elaborate them, because it stops
%       at the end of planning.
%     -> That converts a prohibition into a NAMED model simplification, which is
%        exactly what Transfer rewards. Chapter 2.2 flagged the tension; this
%        paragraph resolves it.
%   - PRODUCT SCOPE DEFINED IN FULL, PROJECT SCOPE ONLY SKETCHED - and why.
%   - THE CHOSEN SPRINT LENGTH AND REFINEMENT CADENCE, and what was traded away.
%     The cadence was derived from PMI's own practice figure and budget ceiling
%     in 3.3 - say so, and say what would change with a more experienced team.
%   - ESTIMATION KEPT TO ONE LOCATING SENTENCE inside sprint planning, with no
%     planning poker. Justify: the report's question is refinement, not effort.
%   - THE FIRST SPRINT DELIBERATELY UNDERLOADED, because a new team judges its
%     own capacity badly.
%   - SCRUM APPLIED AS DESCRIBED, BUT WITH A FICTITIOUS TEAM OF A STATED SIZE -
%     say where reality would differ.
%
% !! DO NOT use the 2004-to-2020 Scrum drift here (visibility -> transparency,
%    ScrumMaster-as-project-manager removed, commitment -> Sprint Goal). It is
%    verified and tempting, but it is SCRUM'S OWN VERSION HISTORY, not an
%    adaptation this project made - it scores nothing under Transfer and would
%    consume the entire citation budget of the chapter.


% --- P3: resources - time, budget, people -------------------------------
% NO CITATION. A named criterion (10 %), and the task PDF asks for it by name:
% "resources such as time, budget and people are used efficiently".
% Two lines of leverage are already in the report and should be collected here,
% in your own words:
%   - the inverted iron triangle from chapter 4: fixing time and cost and
%     letting scope vary is itself a resource decision;
%   - the PO's pricing move from 3.1: "here is the impact if you want this in".
% Also reflect on the report's OWN resource compliance - the page limit is
% explicitly named in the task PDF as a resource requirement.


% --- P4: what would be done differently ---------------------------------
% NO CITATION. [GL] key words to answer: goal achievement, effects and changes
% achieved, areas requiring improvement, continuous improvements to implement,
% how sustainable structures could be created to anchor the results.
% Be concrete and be willing to name a weakness - depth of reflection is what
% the 25 % Quality criterion actually measures.


% --- P5: implications for further study and professional practice -------
% NO CITATION. Answers the [GL] concluding question directly. One short
% paragraph, still in the third person.


% --- P6: closing assessment ---------------------------------------------
% ** THE ONE CLOSING CITATION. ** Pick exactly one of the three below - all
% three are sources already used earlier, so all satisfy the no-new-sources
% rule.
%
%   ***** DECIDED: PMBOK Guide Section 3.2. The two alternatives below are
%   recorded for the audit trail only - do not use them. *****
%   \parencite[PMBOK Guide, Section~3.2]{pmbok2021}
%   "There is no single approach that can be applied to all projects all of the
%    time."
%   -> Fourteen words, PMI's own voice, and it IS the claim rather than a proxy
%      for it. Most economical.
%   !! The part name is mandatory in this locator - the Standard also has a
%      Section 3.2.
%
%   ALTERNATIVE (a) - \parencite[p.~81]{kerzner2022}
%   "Selecting the right framework may seem like a relatively easy thing to do.
%    However, all methodologies and frameworks come with disadvantages as well as
%    advantages."
%   -> The strongest rhetorical choice, because it comes from the most
%      traditional source in the bibliography.
%
%   ALTERNATIVE (b) - \parencite[p.~1520]{mishra2023}
%   "The study came to the conclusion that no single silver belt methodology is
%    appropriate for all kinds of software development projects. The software
%    firms may therefore need to use a hybrid framework that combines agile and
%    waterfall."
%   -> The only peer-reviewed CONCLUSION available, and it names hybrids
%      explicitly.
%   !! "silver belt" is verbatim in the source - the authors write it that way.
%      Reproduce it or paraphrase; do not silently correct it to "silver bullet".
%      The cleaner p. 1518 wording "No methodology is a 'magic solution' for all
%      types of projects" says the same thing without the odd term.
%
% ***** DECIDED: YES, use principle 10 as the second citation. *****
% Chapter 5 allows 0-2. PMBOK Section 3.2 closes the chapter; principle 10 opens
% the Transfer paragraph (P2). Together that is exactly two, and both are
% sources already used earlier in the report, so the no-new-sources rule holds.
% NOTE: this means the locator footnote is needed. The first Beck citation in
% the report is in chapter 4 (principle 2 / value 4), so the footnote goes
% THERE, not here.
%
% SECOND CITATION - the most elegant one available for this paragraph:
%   \parencite[para.~10]{beck2001principles}
%   "Simplicity--the art of maximizing the amount of work not done--is
%    essential."
%   -> Turns the whole P2 list of omissions from a set of shortcuts into a
%      principled position. Moved here from 2.1, where the scope definition won
%      the single available slot.
%   !! The double hyphen is IN THE ORIGINAL, twice. Reproduce it or ellipse
%      around it - do not autocorrect it to an em dash inside quotation marks.
%   !! Needs the locator footnote at the top of this file if no Beck citation
%      appeared in chapter 4.
%
% OWN VOICE, to close: agile fits this product for the reasons given, hybrids
% are legitimate elsewhere, and there is no perfect methodology. Do NOT end on
% "agile is simply better" - a balanced close is both more defensible and better
% aligned with the Quality criterion.
